%%
%% paper77_conway_knoten_de.tex
%% Autor: Michael Fuhrmann
%% Projekt: Specialist Mathematics Project
%% Thema: Conway's Knoten und das Schnittigkeitsproblem in der Knotentheorie
%% Sprache: Deutsch
%% Erstellt: 2026-03-12
%% Letzte Änderung: 2026-03-12
%% Build: 164
%%
%% Beschreibung:
%%   Deutschsprachige Version von Paper 77. Der Conway-Knoten (11 Kreuzungen)
%%   war jahrzehntelang das einzige primäre Knoten mit ≤ 12 Kreuzungen,
%%   dessen glatte Konkordanzklasse unbekannt war. Lisa Piccirillo bewies
%%   2020, dass der Conway-Knoten NICHT glatt scheibig ist. Die topologische
%%   Schnittigkeit bleibt offen.
%%

\documentclass[12pt,a4paper]{amsart}

%% Pakete für Mathematik, Layout, Encoding und Links
\usepackage{amsmath,amssymb,amsthm,mathtools,enumitem,booktabs,array,hyperref}
\usepackage[utf8]{inputenc}
\usepackage[T1]{fontenc}
\usepackage[ngerman]{babel}
\usepackage{geometry}
\geometry{margin=2.5cm}

%% Theorem-Umgebungen (englisch für Nummerierung)
\newtheorem{theorem}{Theorem}[section]
\newtheorem{lemma}[theorem]{Lemma}
\newtheorem{corollary}[theorem]{Korollar}
\newtheorem{proposition}[theorem]{Proposition}

%% Deutsche Theorem-Umgebungen
\newtheorem{satz}[theorem]{Satz}
\newtheorem{korollar}[theorem]{Korollar}
\newtheorem{vermutung}[theorem]{Vermutung}

\theoremstyle{definition}
\newtheorem{definition}[theorem]{Definition}
\newtheorem{definition_de}[theorem]{Definition}
\newtheorem{beispiel}[theorem]{Beispiel}

\theoremstyle{remark}
\newtheorem{bemerkung}[theorem]{Bemerkung}
\newtheorem{remark}[theorem]{Bemerkung}
\newtheorem{conjecture}[theorem]{Vermutung}

%% Mathematische Operatoren
\DeclareMathOperator{\Arf}{Arf}
\DeclareMathOperator{\sign}{sign}
\DeclareMathOperator{\rank}{rank}
\newcommand{\Z}{\mathbb{Z}}
\newcommand{\Q}{\mathbb{Q}}

%% Hyperref-Konfiguration
\hypersetup{
    colorlinks=true,
    linkcolor=blue,
    citecolor=blue,
    urlcolor=blue,
    pdftitle={Conway-Knoten und das Schnittigkeitsproblem in der Knotentheorie},
    pdfauthor={Michael Fuhrmann}
}

%% -------------------------------------------------------------------------
%% Abstract VOR \maketitle
%% -------------------------------------------------------------------------
\begin{abstract}
Der Conway-Knoten $C$, ein Knoten mit 11 Kreuzungen entdeckt von John Horton
Conway, ist das Mutant des Kinoshita-Terasaka-Knotens. Jahrzehntelang war
er der einzige primäre Knoten mit höchstens 12 Kreuzungen, dessen glatte
Konkordanzklasse unbekannt war. Im Jahr 2020 löste Lisa Piccirillo~\cite{Piccirillo2020}
dieses Problem: Sie bewies, dass $C$ \emph{nicht} glatt scheibig ist,
mit einer bemerkenswerten Konstruktion — einem Knoten $K'$ mit demselben
Null-Chirurgie-Ergebnis wie $C$, aber nichtverschwinden\-dem Rasmussen
$s$-Invariant.

Wir präsentieren die Konkordanzgruppe $\mathcal{C}$, die Unterscheidung
zwischen topologischer und glatter Schnittigkeit, die relevanten
Konkordanzinvarianten (Arf-Invariante, Rasmussen $s$-Invariante,
Heegaard-Floer $\tau$-Invariante), Piccirillo's Beweis im Detail und
die verbleibenden offenen Fragen — insbesondere, ob der Conway-Knoten
topologisch scheibig ist.
\end{abstract}

\title{Conway-Knoten und das Schnittigkeitsproblem\\
in der Knotentheorie}
\author{Michael Fuhrmann}
\address{Specialist Mathematics Project, Build 165}
\email{specialist-maths@project.local}
\date{2026-03-12}

\maketitle

%% -------------------------------------------------------------------------
\section{Einleitung}
%% -------------------------------------------------------------------------

Knotenkonkordanz ist eines der zentralen Themen der niedrig-dimensionalen
Topologie. Zwei orientierte Knoten $K_0, K_1 \subset S^3$ sind
\emph{glatt konkordant}, falls eine glatt eingebettete Fläche
$A \cong S^1 \times [0,1]$ in $S^3 \times [0,1]$ existiert mit
$\partial A = K_1 \times \{1\} \cup (-K_0) \times \{0\}$. Konkordanzklassen
bilden eine Gruppe $\mathcal{C}$ unter Zusammenheftung $\#$, die
\emph{glatte Konkordanzgruppe}.

Ein Knoten ist \emph{scheibig} (engl.\ \emph{slice}), wenn er konkordant
zum trivialen Knoten ist, äquivalent wenn er eine glatt eingebettete Scheibe
$D^2 \hookrightarrow B^4$ berandet. Der Unterschied zwischen
\emph{topologisch scheibig} und \emph{glatt scheibig} ist eine der
tiefgreifendsten Manifestationen des Unterschieds zwischen topologischer und
glatter Kategorie in Dimension 4.

Der \emph{Conway-Knoten} $C$ (bezeichnet $11n34$ in der Rolfsen-Tabelle~\cite{Rolfsen1976})
ist das Mutant des Kinoshita-Terasaka-Knotens $KT$ ($11n42$), erhalten durch
Ausschneiden entlang einer Conway-Sphäre und Wiederverkleben mit einer
$180°$-Rotation. Der Conway-Knoten erfüllt alle bekannten notwendigen
Bedingungen für glatte Schnittigkeit: seine Arf-Invariante verschwindet,
sein Alexanderpolynom erfüllt die Fox-Milnor-Faktorisierung, und er ist
algebraisch scheibig. Dennoch blieb seine glatte Konkordanzklasse
jahrzehntelang unbekannt.

In einem wegweisenden Paper~\cite{Piccirillo2020} bewies Lisa Piccirillo (2020):

\begin{satz}[Piccirillo 2020~\cite{Piccirillo2020}]
\label{satz:Piccirillo}
Der Conway-Knoten ist nicht glatt scheibig.
\end{satz}

%% -------------------------------------------------------------------------
\section{Konkordanz und Schnittigkeit}
%% -------------------------------------------------------------------------

\subsection{Die Konkordanzgruppe}

\begin{definition_de}[Knotenkonkordanz]
Zwei orientierte Knoten $K_0, K_1 \subset S^3$ sind \emph{glatt konkordant},
$K_0 \sim K_1$, falls eine glatte eigentliche Einbettung
$f \colon S^1 \times [0,1] \hookrightarrow S^3 \times [0,1]$
existiert mit $f(S^1 \times \{i\}) = K_i$ für $i = 0, 1$.
\end{definition_de}

\begin{satz}[Fox-Milnor~\cite{FoxMilnor1966}]
Die Konkordanzklassen orientierter Knoten in $S^3$ bilden eine abelsche
Gruppe $\mathcal{C}$ unter $\#$, die \emph{glatte Konkordanzgruppe}.
Das neutrale Element ist die Klasse des trivialen Knotens. Das Inverse von
$[K]$ ist $[-K]$ (das Spiegelbild mit umgekehrter Orientierung).
\end{satz}

\subsection{Topologisch vs.\ glatt scheibig}

\begin{definition_de}[Topologisch und glatt scheibig]
Ein Knoten $K \subset S^3$ ist:
\begin{enumerate}[label=\roman*)]
\item \emph{glatt scheibig}, falls er eine glatt eingebettete Scheibe
      $D^2 \hookrightarrow B^4$ berandet;
\item \emph{topologisch scheibig}, falls er eine lokal flache topologisch
      eingebettete Scheibe $D^2 \hookrightarrow B^4$ berandet.
\end{enumerate}
Jeder glatt scheibige Knoten ist topologisch scheibig. Die Umkehrung gilt
nicht: es gibt topologisch, aber nicht glatt scheibige Knoten.
\end{definition_de}

%% -------------------------------------------------------------------------
\section{Der Conway-Knoten und sein Mutant}
%% -------------------------------------------------------------------------

\subsection{Conway-Sphären und Mutation}

\begin{definition_de}[Conway-Sphäre und Mutation]
Eine \emph{Conway-Sphäre} in $S^3 \setminus K$ ist eine eingebettete
$2$-Sphäre, die $K$ in genau 4 Punkten trifft und $K$ in zwei $2$-Tangles
zerlegt. Das \emph{Mutant} $K^m$ von $K$ entsteht durch Ausschneiden entlang
der Conway-Sphäre und Wiederverkleben mittels einer Involution (Rotation um
$180°$ um eine waagrechte oder senkrechte Achse).
\end{definition_de}

\begin{satz}
Folgende Konkordanzinvarianten sind unter Mutation erhalten:
\begin{enumerate}[label=\roman*)]
\item Alexanderpolynom $\Delta_K(t)$;
\item Signatur $\sigma(K)$;
\item Arf-Invariante $\Arf(K)$;
\item Determinante $\det(K) = |\Delta_K(-1)|$.
\end{enumerate}
\end{satz}

Da der Kinoshita-Terasaka-Knoten $KT$ scheibig ist (er berandet eine
Bandscheibe in $B^4$~\cite{KinoshitaTerasaka1957}), verschwinden all diese
Invarianten für $KT$ und damit auch für sein Mutant $C$.

\begin{proposition}
Der Conway-Knoten $C$ ist algebraisch scheibig: Im Witt-Ring von $\Q(t)$
verschwindet die Witt-Klasse seiner Seifert-Paarung.
\end{proposition}

%% -------------------------------------------------------------------------
\section{Konkordanzinvarianten}
%% -------------------------------------------------------------------------

\subsection{Die Arf-Invariante}

\begin{definition_de}[Arf-Invariante]
Die \emph{Arf-Invariante} von $K$ ist $\Arf(K) \in \Z/2\Z$,
berechenbar als $\Delta_K(-1) \pmod{8}/2$ (für ungerades $\Delta_K(-1)$)
oder via die Goeritz-Matrix modulo 8.
\end{definition_de}

\begin{satz}[Robertello 1965~\cite{Robertello1965}]
Wenn $K$ scheibig ist, dann $\Arf(K) = 0$.
\end{satz}

Für den Conway-Knoten: $\Arf(C) = 0$, diese Obstruktion verschwindet also.

\subsection{Das Alexanderpolynom und die Fox-Milnor-Bedingung}

\begin{satz}[Fox-Milnor~\cite{FoxMilnor1966}]
Falls $K$ scheibig ist, gilt $\Delta_K(t) = f(t) \cdot f(t^{-1})$
für ein $f(t) \in \Z[t, t^{-1}]$.
\end{satz}

Das Alexanderpolynom des Conway-Knotens erfüllt diese Bedingung.

\subsection{Die Rasmussen $s$-Invariante}

\begin{definition_de}[$s$-Invariante~\cite{Rasmussen2010}]
Die \emph{Rasmussen $s$-Invariante} ist eine Knotenkonkordanzinvariante
$s \colon \mathcal{K} \to 2\Z$, die aus der Khovanov-Homologie gewonnen wird.
Sie ist ein Gruppenhomomorphismus auf $\mathcal{C}$ und erfüllt:
\begin{enumerate}[label=\roman*)]
\item $s(K) = 0$ falls $K$ scheibig ist;
\item $|s(K)| \leq 2g_4(K)$, wobei $g_4$ das glatte 4-Kugel-Geschlecht ist;
\item $s(T(p,q)) = (p-1)(q-1)$ für Torusknoten.
\end{enumerate}
\end{definition_de}

\begin{satz}[Rasmussen 2010~\cite{Rasmussen2010}]
$s(K)/2$ ist eine untere Schranke für das glatte 4-Kugel-Geschlecht $g_4(K)$:
\[
    |s(K)|/2 \leq g_4(K).
\]
Insbesondere ist $s(K) = 0$ notwendig dafür, dass $K$ glatt scheibig ist.
\end{satz}

Die $s$-Invariante ist im Allgemeinen \emph{nicht} mutationsinvariant.
Piccirillo's entscheidende Beobachtung war, dass der Conway-Knoten einen
„Begleitknoten" mit nichtverschwindender $s$-Invariante besitzt.

\subsection{Die Heegaard-Floer $\tau$-Invariante}

\begin{definition_de}[$\tau$-Invariante~\cite{OzsvathSzabo2003}]
Die Ozsváth-Szabó \emph{$\tau$-Invariante} ist ein Konkordanzhomomorphismus
$\tau \colon \mathcal{C} \to \Z$, definiert via der Knoten-Floer-Homologie
$\widehat{HFK}(K)$. Es gilt:
\begin{enumerate}[label=\roman*)]
\item $\tau(K) = 0$ für scheibige Knoten;
\item $|\tau(K)| \leq g_4(K)$;
\item $\tau(T(p,q)) = (p-1)(q-1)/2$.
\end{enumerate}
\end{definition_de}

\begin{proposition}
$\tau(C) = 0$.
\end{proposition}

\begin{proof}
Dies folgt aus der Mutationsinvarianz von $\tau$ (Ozsváth-Szabó~\cite{OzsvathSzabo2004})
und $\tau(KT) = 0$ (da $KT$ scheibig ist).
\end{proof}

%% -------------------------------------------------------------------------
\section{Piccirillo's Beweis}
%% -------------------------------------------------------------------------

Wir präsentieren die Hauptideen von Piccirillo's Beweis~\cite{Piccirillo2020},
dass der Conway-Knoten nicht glatt scheibig ist.

\subsection{Die Schlüsselkonstruktion}

\begin{definition_de}[Null-Chirurgie]
Die \emph{Null-Chirurgie} $S^3_0(K)$ ist die 3-Mannigfaltigkeit, die aus
$S^3$ durch Dehn-Chirurgie entlang $K$ mit Rahmungskoeffizient 0 entsteht.
\end{definition_de}

\begin{lemma}[Piccirillo~\cite{Piccirillo2020}]
\label{lem:Piccirillo}
Sei $K'$ der Knoten, der aus $C$ durch eine spezifische Modifikation
mittels des Kirby-Diagramms einer exotischen $B^4$ entsteht. Dann:
\begin{enumerate}[label=\roman*)]
\item $S^3_0(K') \cong S^3_0(C)$ (als orientierte 3-Mannigfaltigkeiten);
\item $s(K') \neq 0$ (explizit: $s(K') = -2$).
\end{enumerate}
\end{lemma}

Die Konstruktion von $K'$ nutzt folgende Strategie: Man beginnt mit der
Spur $X_C = B^4 \cup_C h^2$ (die 4-Mannigfaltigkeit mit Rand $S^3_0(C)$,
erhalten durch Ankleben eines 2-Henkels entlang $C$ mit Rahmung 0).
Piccirillo findet einen anderen 2-Henkel $h'^2$ entlang $K'$ mit derselben
Rahmung, sodass $X_{K'} \cong X_C$ (diffeomorphe 4-Mannigfaltigkeiten).
Dies ist ein rein kombinatorisches Kirby-Kalkül-Argument.

\subsection{Freedman's Satz}

\begin{satz}[Freedman 1982~\cite{Freedman1982}]
Ein Knoten $K \subset S^3$ mit verschwindender Arf-Invariante ist
topologisch scheibig. Insbesondere ist jeder Knoten mit trivialem
Alexanderpolynom ($\Delta_K(t) = 1$) topologisch scheibig.
\end{satz}

\subsection{Beweis von Satz~\ref{satz:Piccirillo}}

\begin{proof}[Beweis: $C$ ist nicht glatt scheibig]
\textbf{Schritt 1:} Piccirillo konstruiert den Knoten $K'$ aus
Lemma~\ref{lem:Piccirillo} explizit mittels Kirby-Kalkül auf dem
Kirby-Diagramm von $X_C$.

\textbf{Schritt 2:} Sie berechnet $s(K') = -2 \neq 0$ durch Berechnung
der Khovanov-Homologie von $K'$ (mittels des KnotTheory-Pakets in
\textit{Mathematica}). Dies zeigt insbesondere, dass $K'$ nicht glatt
scheibig ist.

\textbf{Schritt 3:} Angenommen zur Widerlegung, $C$ wäre glatt scheibig
und berandete eine glatte Scheibe $D \subset B^4$.

\textbf{Schritt 4:} Da $S^3_0(K') \cong S^3_0(C)$, gibt es einen
Diffeomorphismus $\phi \colon X_{K'} \to X_C$ der 4-Mannigfaltigkeitsspuren.

\textbf{Schritt 5:} Mit der glatten Scheibe $D$ für $C$ konstruiert man
eine glatte Scheibe für $K'$: Die Spur $X_C$ bettet sich in $B^4$ ein
(da $C$ angenommenermaßen scheibig ist); Zurückziehen via $\phi$ liefert
eine glatte Scheibe für $K'$ in $B^4$.

\textbf{Schritt 6:} Aber $K'$ ist nicht glatt scheibig ($s(K') = -2 \neq 0$).
Widerspruch.

Also ist $C$ nicht glatt scheibig. \qed
\end{proof}

\begin{bemerkung}
Der Beweis verwendet keine tiefe Eichtheorie oder Floer-Homologie — nur:
(1) eine Kirby-Kalkül-Berechnung (Konstruktion von $K'$),
(2) eine kombinatorische Berechnung ($s(K') \neq 0$),
(3) Freedman's topologischen Satz.
Dies macht den Beweis so elegant und unerwartet.
\end{bemerkung}

%% -------------------------------------------------------------------------
\section{Das 4-Kugel-Geschlecht und Bandknoten}
%% -------------------------------------------------------------------------

\begin{definition_de}[4-Kugel-Geschlecht und Bandknoten]
Das \emph{glatte 4-Kugel-Geschlecht} von $K$ ist
\[
    g_4(K) = \min\{ g(F) \mid F \subset B^4 \text{ glatt, } \partial F = K \}.
\]
$K$ ist scheibig genau dann wenn $g_4(K) = 0$.

$K$ ist ein \emph{Bandknoten} (engl.\ \emph{ribbon knot}), falls er eine
Bandscheibe berandet (eine immergierte Scheibe in $S^3$ mit nur
Bandsingularitäten). Jeder Bandknoten ist glatt scheibig.
\end{definition_de}

\begin{vermutung}[Band-Scheiben-Vermutung]
\label{verm:ribbon-slice}
Jeder glatt scheibige Knoten ist ein Bandknoten.
\end{vermutung}

Der Kinoshita-Terasaka-Knoten $KT$ ist ein Bandknoten (und damit scheibig);
der Conway-Knoten $C$ ist nicht scheibig (Piccirillo) und damit kein Bandknoten.

\begin{satz}[Piccirillo 2020]
$g_4(C) \geq 1$.
\end{satz}

Der genaue Wert $g_4(C)$ ist nicht bekannt; es wird vermutet, dass $g_4(C) = 1$.

%% -------------------------------------------------------------------------
\section{Topologische Schnittigkeit des Conway-Knotens}
%% -------------------------------------------------------------------------

Die Frage nach der \emph{topologischen} Schnittigkeit des Conway-Knotens
ist wesentlich subtiler.

\begin{proposition}
Der Conway-Knoten erfüllt alle klassischen notwendigen Bedingungen für
topologische Schnittigkeit:
\begin{enumerate}[label=\roman*)]
\item $\Arf(C) = 0$;
\item $\Delta_C(t) = f(t) \cdot f(t^{-1})$ (Fox-Milnor-Bedingung);
\item $\sigma_\omega(C) = 0$ für alle $\omega \in S^1$ (Levine-Tristram-Signaturen verschwinden);
\item $C$ ist algebraisch scheibig (Witt-Klasse in $W(\Q(t))$ ist null).
\end{enumerate}
\end{proposition}

\begin{vermutung}[Topologische Schnittigkeit des Conway-Knotens]
\label{verm:Conway-top}
Der Conway-Knoten $C$ ist topologisch scheibig.
\end{vermutung}

Diese Vermutung ist (Stand 2026) offen. Das Haupthindernis: Freedman's
Satz gilt unmittelbar nur für Knoten mit trivialem Alexanderpolynom
($\pi_1 \cong \Z$), was für $C$ nicht gilt.

\begin{vermutung}[Topologische vs.\ glatte Konkordanz]
Der Conway-Knoten hat ein nichttriviales Bild in der topologischen
Konkordanzgruppe $\mathcal{C}^{\mathrm{top}}$, d.h.\ er ist topologisch
nicht konkordant zum trivialen Knoten.
\end{vermutung}

Falls diese Vermutung zutrifft, wäre $C$ ein neues Beispiel eines Knotens,
der algebraisch scheibig aussieht, aber nicht einmal topologisch scheibig ist.

%% -------------------------------------------------------------------------
\section{Bandkonkordanz}
%% -------------------------------------------------------------------------

\begin{definition_de}[Bandkonkordanz]
Eine Konkordanz $A \colon K_0 \to K_1$ ist \emph{Band-Konkordanz}, falls
die Projektion $A \subset S^3 \times [0,1] \to [0,1]$ nur Sattelpunkte
als kritische Punkte hat (keine lokalen Maxima).
\end{definition_de}

\begin{satz}[Gordon 1981~\cite{Gordon1981}]
Bandkonkordanz definiert eine Halbordnung auf der Menge der Knoten.
\end{satz}

\begin{vermutung}[Gordon 1981~\cite{Gordon1981}]
\label{verm:Gordon}
Bandkonkordanz ist eine Halbordnung: Falls $K_0$ und $K_1$ in beiden
Richtungen band-konkordant sind, dann $K_0 = K_1$.
\end{vermutung}

Kürzlicher Fortschritt: Agol~\cite{Agol2022} bewies, dass Bandkonkordanz
eine Ordnung auf Homotopieklassen von 3-Mannigfaltigkeiten induziert.

%% -------------------------------------------------------------------------
\section{Weitere Invarianten und Methoden}
%% -------------------------------------------------------------------------

\subsection{Heegaard-Floer-Homologie-Obstruktionen}

Die Knoten-Floer-Homologie $\widehat{HFK}(K)$ liefert die $\tau$-Invariante
und die $d$-Invarianten der Chirurgie-Mannigfaltigkeiten.
Für den Conway-Knoten gilt $\tau(C) = 0$ (rechnerisch verifiziert,
folgt aus Mutationsinvarianz von $\tau$).

\begin{bemerkung}
Die $s$-Invariante ist vermutlich ebenfalls mutationsinvariant, was
$s(C) = s(KT) = 0$ ergeben würde. Dies ist konsistent mit Piccirillo's
Resultat (sie verwendet nur $s(K') \neq 0$ für den Hilfsknoten $K'$,
nicht $s(C)$).
\end{bemerkung}

\subsection{Satelliten-Konkordanzinvarianten}

Für ein Muster $P$ in einem Volltorus und Begleitknoten $K$ gilt für den
Satelliten $P(K)$:
\[
    \tau(P(K)) = w(P) \cdot \tau(K) + \tau(P(U)),
\]
wobei $w(P)$ die Windungszahl ist. Diese Satellitenformeln erlauben die
Gewinnung von Konkordanzinformation aus einfacheren Invarianten.

%% -------------------------------------------------------------------------
\section{Offene Probleme}
%% -------------------------------------------------------------------------

\begin{vermutung}[Topologische Schnittigkeit von $C$]
Der Conway-Knoten ist topologisch scheibig.
\end{vermutung}

\begin{vermutung}[Band-Scheiben-Vermutung, Verm.~\ref{verm:ribbon-slice}]
Jeder glatt scheibige Knoten ist ein Bandknoten.
\end{vermutung}

\begin{vermutung}[Gordon's Vermutung, Verm.~\ref{verm:Gordon}]
Bandkonkordanz ist eine Halbordnung (injektiv bis auf Knotenäquivalenz).
\end{vermutung}

\begin{vermutung}[Glattes 4-Kugel-Geschlecht von $C$]
$g_4(C) = 1$.
\end{vermutung}

\begin{vermutung}[Mutation und $s$-Invariante]
Die Rasmussen $s$-Invariante ist unter Mutation erhalten.
\end{vermutung}

\begin{vermutung}[Glatte Konkordanzgruppe]
$\mathcal{C}$ enthält einen $\Z^\infty$-Summanden und $\Z/2\Z$-Torsion,
aber ihre vollständige Struktur (insbesondere, ob sie ungerade Torsion
enthält) ist unbekannt.
\end{vermutung}

%% -------------------------------------------------------------------------
\section{Zusammenfassung}
%% -------------------------------------------------------------------------

Die Geschichte des Conway-Knotens illustriert, wie die Lücke zwischen
glatter und topologischer Kategorie in Dimension 4 durch konkrete
Knotentheorie erforscht werden kann. Die Auflösung seiner glatten
Schnittigkeit durch Piccirillo 2020 ist ein Triumph cleveren Kirby-Kalküls
kombiniert mit modernen Konkordanzinvarianten. Die zentrale offene Frage
— ob $C$ topologisch scheibig ist — erfordert gänzlich andere Methoden,
da alle klassischen algebraischen Obstruktionen verschwinden und Freedman's
geometrische Techniken nicht direkt anwendbar sind.

\begin{thebibliography}{99}

\bibitem{Piccirillo2020}
L.~Piccirillo,
\textit{The Conway knot is not slice},
Ann.\ of Math.\ \textbf{191} (2020), 581--591.

\bibitem{Conway1970}
J.~H.~Conway,
\textit{An enumeration of knots and links, and some of their algebraic
properties},
in: J.~Leech (Hrsg.), \textit{Computational Problems in Abstract Algebra},
Pergamon Press, Oxford, 1970, 329--358.

\bibitem{KinoshitaTerasaka1957}
S.~Kinoshita, H.~Terasaka,
\textit{On unions of knots},
Osaka Math.\ J.\ \textbf{9} (1957), 131--153.

\bibitem{FoxMilnor1966}
R.~H.~Fox, J.~W.~Milnor,
\textit{Singularities of 2-spheres in 4-space and cobordism of knots},
Osaka J.\ Math.\ \textbf{3} (1966), 257--267.

\bibitem{Rasmussen2010}
J.~Rasmussen,
\textit{Khovanov homology and the slice genus},
Invent.\ Math.\ \textbf{182} (2010), 419--447.

\bibitem{OzsvathSzabo2003}
P.~Ozsváth, Z.~Szabó,
\textit{Knot Floer homology and the four-ball genus},
Geom.\ Topol.\ \textbf{7} (2003), 615--639.

\bibitem{OzsvathSzabo2004}
P.~Ozsváth, Z.~Szabó,
\textit{Knot Floer homology and integer surgeries},
Algebr.\ Geom.\ Topol.\ \textbf{8} (2008), 101--153.

\bibitem{Freedman1982}
M.~H.~Freedman,
\textit{The topology of four-dimensional manifolds},
J.\ Differential Geom.\ \textbf{17} (1982), 357--453.

\bibitem{KronheimerMrowka1993}
P.~B.~Kronheimer, T.~S.~Mrowka,
\textit{Gauge theory for embedded surfaces, I},
Topology \textbf{32} (1993), 773--826.

\bibitem{Levine1969}
J.~P.~Levine,
\textit{Knot cobordism groups in codimension two},
Comment.\ Math.\ Helv.\ \textbf{44} (1969), 229--244.

\bibitem{Robertello1965}
R.~A.~Robertello,
\textit{An invariant of knot cobordism},
Comm.\ Pure Appl.\ Math.\ \textbf{18} (1965), 543--555.

\bibitem{Gordon1981}
C.~McA.~Gordon,
\textit{Some aspects of classical knot theory},
in: \textit{Knot Theory}, Lecture Notes in Math.\ \textbf{685},
Springer, Berlin, 1978, 1--60.

\bibitem{Agol2022}
I.~Agol,
\textit{Ribbon concordance of knots is a partial order},
Ann.\ of Math.\ \textbf{196} (2022), 675--698.

\bibitem{Hedden2009}
M.~Hedden,
\textit{Notions of positivity and the Ozsváth-Szabó concordance invariant},
J.\ Knot Theory Ramifications \textbf{19} (2010), 617--629.

\bibitem{Rolfsen1976}
D.~Rolfsen,
\textit{Knots and Links},
Publish or Perish, Berkeley, 1976.

\end{thebibliography}

\end{document}
